\section{Basis}

\subsection{Stockfish}
Wir möchten unsere Engine mit Stockfish vergleichen. Nicht, weil wir davon ausgehen, dass unsere Engine mitthalten kann,
sondern vielmehr, um daraus zu lernen und einen Referenzpunkt zu haben.

Die grundlegende Idee besteht darin, \gls{glos:JUnit}-Tests zu schreiben, welche mit Stockfish interagieren. Dafür muss
Stockfish von der Maschine aus erreichbar sein, auf welcher der \gls{glos:JUnit}-Test ausgeführt wird.

Hierfür haben wir \href{https://testcontainers.com/}{Testcontainers} verwendet. Damit ist es möglich, während der
Testausführung beliebige Docker-Container zu starten und diese nach Abschluss der Tests wieder zu stoppen.

Da Stockfish selbst weder ein Docker-Image noch ein Dockerfile veröffentlicht, haben wir mithilfe von
(\cite{stockfishCompilingFromSource}) und (\cite{socatExamples}) ein eigenes Dockerfile erstellt. Daraus können wir
ein Image bauen und einen Container starten.

Um während der Tests mit Stockfish kommunizieren zu können, benötigen wir einen Client. Ein weit verbreitetes
Kommunikationsprotokoll, das von vielen Schachengines und auch von Stockfish verwendet wird, ist das
\href{https://backscattering.de/chess/uci/}{UCI-Protokoll}. Damit kann Stockfish beispielsweise eine Position übermittelt
und der nächste Zug berechnet werden. Die von Stockfish unterstützten UCI-Befehle sind
\href{https://official-stockfish.github.io/docs/stockfish-wiki/UCI-&-Commands.html}{hier} aufgeführt.

Sämtliche Klassen und Dateien, die entwickelt wurden, um Stockfish während der Tests auszuführen, befinden sich unter
\textit{src/test/java/ch/hslu/cas/msed/blobfish/stockfish} sowie \textit{src/test/resources/stockfish}.

\subsection{\gls{glos:JUnit}-Extension}

Um möglichst einfach innherhalb einer Testmethode mit Stockfish kommunizieren zu können, haben wir uns entschieden,
eine \gls{glos:JUnit}-Extension zu schreiben. Dies ermöglich uns beispielsweise über Callbacks, Code vor oder nach einer
Testklasse oder Testmethode auszuführen (\cite{junitCallbacks}). Ebenfalls kann über PostProcessing Objekte in Felder der
Testklasse injected werden (\cite{junitExtensionRegistration}).

So sehen beispielsweise die ersten Zeilen der Testklasse \textit{QualityTest}, welche eigene Extension verwendet:

\lstinputlisting[
    firstline=32,
    lastline=39,
    firstnumber=1,
    caption={QualityTest},
    label={lst:qualitTestHead}
]{../src/test/java/ch/hslu/cas/msed/blobfish/QualityTest.java}

Mittels der Annotation \textit{@ExtendWith(StockfishExtension.class)} wird die eigene Extension geladen. Dies bewirkt,
dass bevor die erste Testmethode ausgeführt wird, kontrolliert wird, ob lokal eine Dockercontainer von Stockfish bereits
gestartet ist. Falls dies nicht der Fall ist, wird ein neuer Container gestartet.
Als nächstes, wird eine der StockFishService-Klasse, eine Fassade, innerhalb der Extension initialisiert. Diese Objekt
wird dann mittels PostProcessing in das Feld auf Zeile 3 injected.

Annschliessend werden die Testmethoden ausgeführt. Innerhalb der Testmethoden kann nun via der \textit{stockFishService}
Variable mit Stockfish kommuniziert werden, und so beispielsweise Positionen berechnet werden.
Sicherheitshalber wird die Extension zwischen den Testmethoden der Befehl \href{https://official-stockfish.github.io/docs/stockfish-wiki/UCI-\&-Commands.html\#ucinewgame}{ucinewgame}
an Stockfish übermittelt, um Stockfish zu resetten.

Nachdem alle Testmethoden ausgeführt sind und ein Dockercontainer für die Testausführung gestartet wurde, stoppt dieser
die Extension.